\section{Results}
\label{sec:results}

\subsection{The check-in-level representation improves the category task}
\label{sec:results-part1}
The check-in-level representation reaches a higher next-category macro-F1 than a standard place
embedding (HGI~\cite{huang2023hgi}) at every dataset
(Table~\ref{tab:substrate}). Only the input representation differs, a vector per visit versus a vector per place. The check-in-level column is the dedicated category model of Table~\ref{tab:results} at seed 0, and the place-level column swaps its input.
The gaps run from $+0.23$ to $+6.29$ points and stay under one point at the three largest
datasets, so what the comparison establishes is a consistent direction rather than a large effect.

The benefit is category-only by construction: the region task reads region-level vectors, not the
per-visit vectors (Section~\ref{sec:method-model}).
CTLE~\cite{lin2021ctle}, the closest prior contextual embedding, also gives each visit its own vector, but it pretrains on place identifiers and timestamps alone, so the category vocabulary never enters its training, whereas our graph reads each visit's category and time of day as input features. At Florida, we fine-tune CTLE together with the task model, using its authors' defaults and the same 64 dimensions and windowing as ours. It reaches $33.45$ macro-F1 at its best epoch and $29.69$ at its final one (an epoch is one pass over the training data), below the place embedding under the same best-epoch rule.
With frozen weights, fed to the same single-task model, CTLE repeats the ordering at Alabama, Arizona, and Istanbul. These CTLE results establish the ordering between the two representation families rather than a
difference against the values in Table~\ref{tab:substrate}.
A feature-concatenation control joins the place embedding with the same raw
per-visit features that our graph reads (the category one-hot plus hour and
day-of-week encodings) and feeds the result to the same single-task model. On
the category axis, at Alabama, Arizona, and Florida this control raises the
place embedding by $+2.0$, $+1.7$, and $+0.8$ macro-F1, against
place-to-check-in gaps of $+1.62$, $+2.58$, and $+0.23$ at the same states
(Table~\ref{tab:substrate}). Most of the category difference is
therefore already available in the raw per-visit features, and this control does not separate
what the check-in-level representation adds beyond them on this axis.




\subsection{One model, two tasks}
\label{sec:results-part2}
One model serves both tasks, and every difference on both axes stays inside the bound its axis
supports. On next region it
outperforms the dedicated model at Texas and California, the two datasets with the most regions, and
at the other four it stays within the two-point margin registered before any result was read
(statistical non-inferiority, assessed with two one-sided tests). On next category it outperforms the
dedicated model at Florida, and at every dataset the difference between the two is equivalent to zero
within half a point, simultaneously across all six.

Half a point is worth placing against the two comparisons that give the category axis its scale. The
input representation moves the same metric by $+0.23$ to $+6.29$ points, up to thirty-two times the
largest difference the choice of architecture makes, and the joint model stands at least $3.06$ points
above the strongest external baseline at every dataset. Whether one model or two produce the category
prediction is therefore the smallest of the three effects at work here.
For scale: always predicting the most common category reaches between $5.7$ and $7.3$ macro-F1
depending on the dataset (the majority-class floor, lowest at Florida, whose category mix is the
most even), and a random region top-ten guess is right at most about two percent of the time
(Section~\ref{sec:setup-metrics}).

Every reported model is one saved model per fold, read at its validation-selected epoch: each
dedicated model at its task's best epoch, and the joint model at the epoch selected by the geometric
mean of its two validation metrics, with both tasks read from that one model, which is what a
deployed system can serve. Reading each task at its own best epoch instead would describe no single
saved model, and it is more favorable to the joint model, by at most $0.23$ macro-F1 and $0.93$
Acc@10 points at any one seed, enough to turn four further category comparisons and two further
region comparisons into improvements that survive the same Holm correction. Every verdict in this
paper is the stricter convention's.


On region (Acc@10), the joint model outperforms the dedicated model at Texas and California. The two
datasets where it outperforms are the two with the largest region vocabularies, which is where the
region task has the most classes to rank.
The ordering is not monotone inside that pair, since California has more regions than Texas and a
slightly smaller gain, and region count co-varies with corpus size across these states, so the
grouping is read as an observation about where the benefit appears rather than as a law relating
gain to region count. Why the benefit appears there is a separate question (Section~\ref{sec:discussion}).

On category, Florida is the one dataset where the difference survives the
Holm correction across the six ($+0.19$; 90\% CI $+0.14$ to $+0.25$; corrected $p=0.011$), with 19 of the
20 folds favoring the joint model. The other five differences do not survive it, and because the
equivalence margin was registered for the region axis only, they are reported by the bound their intervals support
rather than as matches: Istanbul ($+0.08$; $+0.01$ to $+0.15$), Arizona ($-0.00$; $-0.04$ to
$+0.03$), California ($-0.00$; $-0.03$ to $+0.02$), Texas ($-0.13$; $-0.19$ to $-0.08$) and Alabama
($-0.19$; $-0.33$ to $-0.04$).
The intervals themselves carry what can be said about magnitude:
the widest of them reaches $0.33$ points from zero, at Alabama, and with a Bonferroni adjustment
across the six comparisons the widest reaches $0.49$, so all six differences stay within half a point
of zero at once.
The bound is read off the intervals rather than established by a
further test.

On region, Texas ($+1.21$; $+1.13$ to $+1.29$; corrected $p=0.00013$) and California ($+1.06$;
$+1.03$ to $+1.08$; corrected $p<10^{-4}$) survive the correction, with all 20 folds favoring the
joint model at both. The remaining four stay within the registered two-point margin, and their
direction is stated rather than rounded away: all four are deficits, at Alabama ($-0.87$; $-1.00$
to $-0.75$), Arizona ($-0.44$; $-0.62$ to $-0.25$), Florida ($-0.16$; $-0.19$ to $-0.13$) and
Istanbul ($-0.08$; $-0.16$ to $-0.002$), and every one of the four intervals lies entirely below
zero. A test in the reverse direction, applied post hoc to the same six region comparisons and
corrected across them, resolves three of the four, and Istanbul's is not resolved. Each of the four
stays inside the margin, but none of them is a tie.


The joint model is also above every external baseline reported, on both tasks,
across all six datasets (Table~\ref{tab:results}): on region, the primary
region-native comparison HMT-GRN~\cite{lim2022hmtgrn}, STAN~\cite{luo2021stan}
trained from the raw check-ins, and a ReHDM reference~\cite{li2025rehdm}; on
category, POI-RGNN~\cite{capanema2023poirgnn} and the Markov-K floor. These externals run on their own embeddings, so this comparison
also includes the representation advantage of Section~\ref{sec:results-part1}; the
controlled comparison for the joint model remains the Dedicated column of
Table~\ref{tab:results}. The first-order Markov region floor of Section~\ref{sec:setup-baselines},
computed under our windows and folds, reaches $51$ to $72$ Acc@10 across the
datasets; the joint model exceeds it by $4.1$ to $10.0$ points on all six
datasets.

That floor is also above the three external region systems at most datasets: HMT-GRN falls below
it at all six, the ReHDM reference at three, and STAN at four. The floor reads a strong signal that
our sliding windows expose, since they advance one visit at a time and the target region is the
last visited region in $32.9$ percent of the windows at Alabama, and the three systems do not meet
it on equal terms (HMT-GRN on our data, folds, and initialization, STAN on our folds with its own
embeddings and sequences, ReHDM under its own published protocol). We therefore treat the floor,
not the external systems, as the reference the region task has to clear, and the dedicated and
joint models are above it at every dataset.

At Istanbul, the non-U.S. city, both differences are within a tenth of a point of the dedicated
models (category $+0.08$ macro-F1, region $-0.08$ Acc@10), so what repeats on a different continent
and region unit is the bounded picture, with neither difference resolved as a gain.
